\needspace{12\baselineskip}
\SourceLeaf{319}
\pdfbookmark[1]{16. Über die Zodiaken der Himmelskörper}{zusatz-016}
\begin{center}
[Über die Zodiaken der Himmelskörper.]
\end{center}

\begin{center}
[16.]
\end{center}

\par\noindent\rule{\textwidth}{.3pt}\par

\noindent\small Astronomische Nachrichten. Band XXVII. Nr. 625. Seite 1. 1848 Januar 29.\normalsize

\par\noindent\rule{\textwidth}{.3pt}\par

\begin{center}
Schreiben des Herrn Geheimen Hofraths \textsc{Gauss} an den Herausgeber.

Göttingen 1847 November 23.
\end{center}

Da die in Nr. 615 der A.\,N. abgedruckten zweiten Elemente der Iris, welche Herr Prof. \textsc{Goldschmidt} gleich nach dem 19\textsuperscript{ten} September berechnet hatte, auch noch in den beiden folgenden Monaten eine gute Übereinstimmung mit den Beobachtungen gezeigt haben (am 17\textsuperscript{ten} und 18\textsuperscript{ten} November war die Differenz von Herrn \textsc{Rümkers} Meridianbeobachtung \(16^{\prime\prime}\) in ger. Aufst.\ und \(2^{\prime\prime}\) in der Abweichung), so ist man berechtigt, sie als schon sehr genähert zu betrachten, und ich habe deshalb Herrn Prof. \textsc{Goldschmidt} veranlasst, danach den Zodiakus der Iris zu berechnen. Das Resultat dieser Arbeit, welches zur Erleichterung der Nachforschungen auf Identität in frühern Beobachtungen wird dienen können, lasse ich hier nachfolgen.

Indem ich bei dieser Veranlassung den Aufsatz\HistoricalNoteEight{Band VI, S. 106.]} wieder in die Hände nehme, in welchem ich vor 44 Jahren die allgemeine Methode zur Bestimmung der Limiten eines solchen Zodiakus gegeben habe, sehe ich, dass ich darin auch schon die Ausnahmefälle angedeutet habe, wo das Feld der geocentrischen Erscheinung eines die Sonne nach \textsc{Kepler}schen Gesetzen umkreisenden Himmelskörpers auf der Himmelskugel entweder gar keine oder nur Eine Limite hat, obgleich die Methode im erstern Falle eine in sich zurücklaufende
\SourceLeaf{320}Linie, im andern zwei solche Linien ergibt. Auch ist die Frage daselbst
aufgeworfen, was denn in solchen Ausnahmsfällen diese durch Rechnung
gefundenen Linien eigentlich bedeuten. Ich habe mich damals auf diese
Andeutungen beschränkt, weil eine weitere Ausführung dort ein Horsd’oeuvre
gewesen wäre, und ich auch gern andern das Vergnügen lassen wollte, sich mit
einer meiner Meinung nach nicht uninteressanten mathematischen Aufgabe zu
beschäftigen. Da mir jedoch nicht bekannt geworden ist, dass ein anderer in
der langen Zwischenzeit die mir inzwischen ganz aus dem Gedächtniss gekommene
Frage aufgenommen hätte, so ergreife ich diese Gelegenheit, um wenigstens den
Hauptnerv des zur Beantwortung Nöthigen hier mitzutheilen.

Ein geocentrischer Ort des in Rede stehenden Planeten (oder Cometen) geht
hervor, indem man einen Punkt der Bahn des letztern mit einem Punkt der
Erdbahn combinirt; es kann aber auch einerlei geocentrischer Ort aus zwei,
drei oder vier verschiedenen Combinationen hervorgehen. Um die Vorstellungen
zu fixiren, lege man durch die Sonne eine Ebene, gegen welche die einen
vorgegebenen geocentrischen Ort, auf der Himmelskugel, repräsentirende gerade
Linie normal ist, und projicire orthographisch auf diese Ebene sowohl die
Erdbahn als die Planetenbahn. Beide Projectionen sind Ellipsen, oder
allgemeiner Kegelschnitte, die sich entweder garnicht, oder zweimal oder
viermal schneiden werden; eine Berührung ist dabei wie das Verschmelzen
zweier Schnitte zu betrachten. Durch jeden Schnitt wird ein geocentrischer
Ort vorgestellt, der entweder mit dem vorgegebenen identisch oder ihm auf der
Himmelskugel entgegengesetzt ist, je nachdem von den beiden Punkten der
Planetenbahn und der Erdbahn, deren Projection zusammenfällt, der erstere oder
der letztere höherliegend ist; als obere Seite der Projectionsebene diejenige
betrachtet, auf welcher der vorgegebene geocentrische Ort liegt.

Es ist hieraus klar, dass wenn irgend ein Punkt der Himmelskugel, fragweise
als geocentrischer Ort, aufgegeben wird, er dies entweder auf gar keine Weise
oder auf Eine Art, oder auf zwei, drei oder vier Arten sein kann. Für einen
gegebenen Planeten scheidet sich so die ganze Fläche der Himmelskugel in
verschiedene Theile, und die nach der von mir gegebenen Methode bestimmten
Linien sind, allgemein zu reden, nichts anderes, als Scheidungen zwischen
zwei Flächentheilen der Kugel, wo die auf der einen Seite liegenden Punkte
auf zwei Arten mehr geocentrische Örter sein können als die auf der
\SourceLeaf{321}andern. Die in den Scheidungslinien liegenden Punkte machen den Übergang,
d. i. sie sind auf Eine Art mehr als die Punkte auf der einen Seite, und auf
Eine Art weniger als die Punkte auf der andern Seite fähig geocentrische Örter
zu sein. Übrigens lässt sich auch das Criterium angeben, wonach a priori
entschieden wird, auf welcher Seite der Scheidungslinie zwei Auflösungen mehr
stattfinden als auf der andern, wobei ich jedoch gegenwärtig mich nicht
aufhalten will\HistoricalNoteEight{Siehe Notiz 17.]}.

Alle bisher bekannten Planeten haben solche Bahnen, dass die durch die Theorie
gefundenen Scheidungslinien immer wahre Limiten sind. Es sind nemlich zwei
Scheidungslinien, die die Himmelskugel in drei Flächenräume abtheilen; zwei
sind isolirte Flächen, in welche gar keine geocentrische Örter fallen, der
dritte zwischen jenen, gürtelartig, enthält alle geocentrischen Örter und jeder
Punkt innerhalb des Gürtels kann auf zwei Arten, jeder Punkt auf der Limite
auf Eine Art geocentrischer Ort sein.

Es lassen sich aber fingirte Bahnen denken\HistoricalNoteDouble{Siehe Notiz 18.]} (Cometenbahnen werden
vielleicht mehrere in dem Fall sein, worüber ich jedoch bisher keine
Nachforschung gemacht habe), wo zwei Limiten die Himmelskugel auch in drei
Stücke scheiden, und wo der eine Theil gar keine geocentrische Örter enthält,
der andere die geocentrischen Örter zu zwei Arten, der dritte hingegen die
Punkte, die auf vier verschiedene Arten geocentrische Örter sein können.

Noch weitere Mannigfaltigkeit ergibt sich, wenn Eine Scheidungslinie sich
selbst einmal oder zweimal schneidet, eine einfache oder doppelte Schleife
bildet. Im ersten Fall wird sie zwei, im zweiten drei Flächenräume von dem
gürtelförmigen Theile abtrennen, in denen resp. 4 und 0 oder 4, 0, 4
Auflösungsarten stattfinden.

Noch anders verhält es sich mit einer Bahn, die in die Erdbahn wie ein
Kettenring eingreift. In einem solchen Fall ist nur Eine zusammenhängende
Limitenlinie, also zwei geschiedene Flächentheile, wo die Punkte des einen
Theils auf Eine Art, die des andern Theils auf drei Arten geocentrische Örter
sein können.

\needspace{29\baselineskip}
\SourceLeaf{322}
\begin{center}
Zodiakus der Iris.
\end{center}
\begingroup\small
\setlength{\tabcolsep}{7pt}
\renewcommand{\arraystretch}{1.15}
\begin{center}
\begin{tabular}{|r|r@{\;}r|r@{\;}r|}
\multicolumn{1}{|c|}{Gerade}&\multicolumn{4}{c|}{Abweichung}\\
\multicolumn{1}{|c|}{Aufst.}&\multicolumn{2}{c|}{nördl. Grenze}&\multicolumn{2}{c|}{südl. Grenze}\\
\(0^{0}\)&\(+12^{0}\)&\(9^{\prime}\)&\(+3^{0}\)&\(44^{\prime}\)\\
5&14&10&5&41\\
10&16&5&7&33\\
15&17&51&9&19\\
20&19&30&10&57\\
25&21&0&12&29\\
30&22&21&13&52\\
35&23&33&15&7\\
40&24&36&16&14\\
45&25&29&17&11\\
50&26&12&17&59\\
55&26&45&18&39\\
60&27&9&19&8\\
65&27&24&19&29\\
70&27&29&19&41\\
75&27&24&19&42\\
80&27&11&19&36\\
85&26&48&19&19\\
90&26&15&18&54\\
95&25&33&18&19\\
100&24&41&17&35\\
105&23&41&16&42\\
110&22&32&15&41\\
115&21&14&14&31\\
120&19&47&13&12\\\hline
\end{tabular}
\end{center}
\begin{center}
\begin{tabular}{|r|r@{\;}r|r@{\;}r|}
\multicolumn{1}{|c|}{Gerade}&\multicolumn{4}{c|}{Abweichung}\\
\multicolumn{1}{|c|}{Aufst.}&\multicolumn{2}{c|}{nördl. Grenze}&\multicolumn{2}{c|}{südl. Grenze}\\
\(125^{0}\)&\(+18^{0}\)&\(12^{\prime}\)&\(+11^{0}\)&\(46^{\prime}\)\\
130&16&29&10&12\\
135&14&40&8&31\\
140&12&45&6&45\\
145&10&44&4&53\\
150&8&39&2&56\\
155&6&31&\(+0\)&56\\
160&4&22&\(-1\)&6\\
165&\(+2\)&12&3&8\\
170&0&0&5&11\\
175&\(-2\)&9&7&13\\
180&4&14&9&11\\
185&6&16&11&6\\
190&8&12&12&56\\
195&10&2&14&40\\
200&11&45&16&19\\
205&13&20&17&50\\
210&14&46&19&13\\
215&16&4&20&58\\
220&17&13&21&35\\
225&18&12&22&33\\
230&19&1&23&23\\
235&19&40&24&4\\
240&20&10&24&35\\\hline
\end{tabular}
\end{center}
\begin{center}
\begin{tabular}{|r|r@{\;}r|r@{\;}r|}
\multicolumn{1}{|c|}{Gerade}&\multicolumn{4}{c|}{Abweichung}\\
\multicolumn{1}{|c|}{Aufst.}&\multicolumn{2}{c|}{nördl. Grenze}&\multicolumn{2}{c|}{südl. Grenze}\\
\(245^{0}\)&\(-20^{0}\)&\(29^{\prime}\)&\(-24^{0}\)&\(58^{\prime}\)\\
250&20&38&25&11\\
255&20&38&25&15\\
260&20&27&25&11\\
265&20&5&24&56\\
270&19&34&24&33\\
275&18&52&24&1\\
280&18&0&23&19\\
285&16&58&22&28\\
290&15&46&21&28\\
295&14&25&20&20\\
300&12&54&19&2\\
305&11&14&17&36\\
310&9&26&16&2\\
315&7&30&14&21\\
320&5&28&12&34\\
325&3&21&10&40\\
330&\(-1\)&9&8&41\\
335&\(+1\)&6&6&39\\
340&3&22&4&34\\
345&5&37&2&28\\
350&7&51&\(-0\)&22\\
355&10&2&\(+1\)&42\\
360&12&9&3&44\\\hline
\end{tabular}
\end{center}
\endgroup
\par\medskip\begin{center}\rule{.22\textwidth}{.3pt}\end{center}
